\section{Discussion}

% \noindent\textbf{Ethics.} The three bugs we exploited were from syzbot's historical reports. All three bugs already have been patched in both upstream and major downstream kernels prior to our analysis. All exploit development was conducted in controlled environments for academic purposes only.

%These bugs were selected from syzbot's historical reports, and no responsible disclosure was necessary.

\noindent\textbf{Portability of Pointer Materialization.} Modern exploitation techniques weaken the assumption that KASLR and SMAP alone prevent attackers from constructing valid kernel pointers. SyzPass therefore models both pointer materialization and object matching. Pointer materialization captures stronger post-disclosure pointer capabilities, while object matching provides a portable strategy when compatible spray objects exist across deployment environments.
% \zhiyun{Discuss how we may improve the scalability of symbolic execution and reach the end of syscalls for more bugs, e.g., using a more performant symbolic execution engine such as qemu-based, e.g., SymFit, or instrumentation-based, SymCC.}
% The speed of symboilc execution becomes a bottleneck of this work and potential hinders more effective effects and path discovery. The efficiency of Angr + memory snapshot based symbolic execution is inherently slow. Thus, the speed of symbolic execution may become a bottleneck of our analysis

\noindent\textbf{Implications of Different Side Effects.} Different side effects have varied implications on exploit development due to their potential to destabilize the system or disrupt the attacker's control. In this work, neutralization focuses on lethal pointer-shaped residual effects, especially dereferences of invalid pointers, since such operations frequently lead to immediate system crashes. UAF write side effects can introduce second-order object-state questions; SyzPass treats them through primitive-aware preservation rules and focuses neutralization on pointer-shaped residual effects.
% the attacker's control over the kernel state
%unusable or causing a crash when the kernel accesses the corrupted field. While sometimes tolerable, such corruption can destroy carefully crafted heap layouts or trigger dereferences of invalid pointers, undermining the exploit.
%Dereferences of invalid pointers are especially problematic, as they are prone to causing system crashes.
%Dereferencing invalid read or write pointer can lead to kernel oops that kills the current process and even kernel panic (if panic\_on\_oops is on); Dereferencing invalid code pointer typically leads to kernel panic.
% Therefore, understanding the nature of side effects is critical: some (e.g., simple overwrites) may be manageable or neutralizable, while others (e.g., dereferencing corrupted pointers) demand strict avoidance or precise neutralization using carefully selected heap objects.
% This variability in consequences motivates the need for fine-grained per-path analysis, as automated by SyzPass, to ensure chosen primitives are paired with side-effect-free or side-effect-resolved paths for stable exploit construction.

\noindent\textbf{Coverage and Runtime Scope.} SyzPass uses time-bounded symbolic execution to keep per-bug primitive analysis tractable. The reported counts summarize the explored primitive-aware paths in the studied cases. Branch-set collection and dereference-depth pruning are engineering mechanisms that make the guidance pipeline practical, while faster symbolic execution backends, such as QEMU-based engines~\cite{symfit} or instrumentation-based approaches~\cite{symcc}, can further expand coverage.

\noindent\textbf{Suppressing Side Effects via Kernel Stalling.} Beyond avoidance and neutralization, attackers may suppress side effects by stalling kernel execution. For example, leveraging FUSE to block \texttt{copy\_from\_user()} can hang the kernel thread during user-kernel data transfer~\cite{fuze_exp}. This prevents subsequent side-effect-inducing instructions from executing and complements SyzPass when a suitable blocking point exists between the primitive and the side effects.

%While this does not yield a complete exploit path on its own, it can serve as a useful workaround to suppress crashes in cases where clean avoidance or neutralization is infeasible, and can be especially valuable in information disclosure or staging multi-step exploitation.

%\noindent\textbf{Non-Deterministic UAF Timing and Side Effects.} In certain UAF vulnerabilities~\cite{CVE20230266}, the timing of the free operation is non-deterministic relative to a series of object uses. That is, the vulnerable object may be accessed multiple times, and the actual free may occur at any point along these accesses depending on scheduling, race conditions, or subtle memory management behavior. As a result, the triggered side effects can vary across different runs of the same exploit—though along the same execution path. For example, if the object is freed after a use, the use is a benign memory access; if freed before, the use could is a side effect or a primitive. This introduces non-determinism in side effect manifestation, complicating both exploit development and symbolic execution, as the set of effects cannot be treated as static. A more effective analysis must therefore consider multiple interleavings of use and free events. % to identify robust, side-effect-resilient exploitation paths

% For peer review papers, you can put extra information on the cover
% page as needed:
% \ifCLASSOPTIONpeerreview
% \begin{center} \bfseries EDICS Category: 3-BBND \end{center}
% \fi
%
% For peerreview papers, this IEEEtran command inserts a page break and
% creates the second title. It will be ignored for other modes.
% \IEEEpeerreviewmaketitle
